% !TEX program = xelatex
% 中文学术译文模板 —— 用 xelatex 编译（ctex 负责中文排版与字体回退）。
% 正文单栏、便于阅读；公式用 LaTeX 重排；图片嵌入 figures/ 下提取的原图。
\documentclass[11pt]{article}

\usepackage[UTF8, scheme=plain]{ctex}        % 中文支持；scheme=plain 不改西文版式
\usepackage[a4paper, margin=2.5cm]{geometry}
\usepackage{amsmath, amssymb, amsfonts}      % 数学公式
\usepackage{graphicx}                         % 插图
\usepackage{booktabs}                         % 三线表
\usepackage{array, multirow, makecell}        % 复杂表格
\usepackage{caption}
\usepackage{float}                            % [H] 精确定位浮动体
\usepackage{url}
\usepackage{enumitem}
\usepackage{xcolor}
\definecolor{linkblue}{RGB}{40,80,150}    % 引用/交叉引用链接色
\definecolor{urlblue}{RGB}{30,90,140}
% hyperref 放在大多数宏包之后；colorlinks 让正文引用、交叉引用、目录可点击跳转
\usepackage[colorlinks=true, breaklinks=true]{hyperref}
\hypersetup{
  citecolor=linkblue,   % 正文 [n] 引用标记
  linkcolor=linkblue,   % 图/表/公式/章节交叉引用
  urlcolor=urlblue,
}
\usepackage{cleveref}                     % \cref 生成"图~1""式~(3)"等带链接的引用
% cleveref 默认英文名，改成中文图表名
\crefname{figure}{图}{图}
\crefname{table}{表}{表}
\crefname{equation}{式}{式}
\crefname{section}{第}{第}
\creflabelformat{section}{#2#1#3 节}      % \cref{sec:x} -> "第 2 节"

% 图表标题中文化：图 1 / 表 1
\captionsetup{labelsep=quad}
\renewcommand{\figurename}{图}
\renewcommand{\tablename}{表}
\graphicspath{{figures/}{work/figures/}{./}}

% 行距略松，学术阅读更舒适
\linespread{1.25}

\title{\bfseries 论文中文标题（括注原英文标题）}
\author{原作者姓名（保留原文）}
\date{}

\begin{document}
\maketitle

\begin{abstract}
\noindent 摘要正文的中文译文……
\end{abstract}

\section{引言}
正文译文。首次出现的术语括注英文，例如卷积神经网络（Convolutional Neural
Network, CNN）。正文引用用 \verb|\cite|，会生成可点击跳转到参考文献的链接，
如已有方法~\cite{ref1} 与~\cite{ref2}。如~\cref{fig:overview} 所示，
本文方法的结果见~\cref{tab:results}。行内公式如 $E=mc^2$，独立公式见~\cref{eq:loss}：
\begin{equation}
  \mathcal{L} = -\sum_{i} y_i \log \hat{y}_i .
  \label{eq:loss}
\end{equation}

% —— 嵌入提取出的原图，图注译为中文 ——
\begin{figure}[H]
  \centering
  \includegraphics[width=0.8\linewidth]{p03_fig1.png}
  \caption{中文图注译文（必要时保留坐标轴等原文标签）。}
  \label{fig:overview}
\end{figure}

% —— 表格重排为中文三线表 ——
\begin{table}[H]
  \centering
  \caption{中文表注译文。}
  \label{tab:results}
  \begin{tabular}{lcc}
    \toprule
    方法 & 准确率 & 参数量 \\
    \midrule
    基线 & 0.91 & 5M \\
    本文 & \textbf{0.95} & 4M \\
    \bottomrule
  \end{tabular}
\end{table}

% 参考文献保留原文，不翻译。用 thebibliography + \bibitem，配合正文 \cite，
% hyperref 会自动把正文 [n] 变成可点击、跳转到对应条目的超链接。
% \bibitem 的标签（ref1、ref2…）与正文 \cite 一一对应，照原文顺序逐条录入。
\begin{thebibliography}{99}
  \bibitem{ref1} Author A, Author B. Title of the first paper. \textit{Venue}, Year.
  \bibitem{ref2} Author C, Author D. Title of the second paper. \textit{Venue}, Year.
\end{thebibliography}

\end{document}
